\documentclass[12pt]{article}  % Set document class and font size

\usepackage{graphicx}           % For including graphics
\usepackage{amsmath}            % For mathematical symbols
\usepackage{amsfonts}           % For mathematical fonts
\usepackage{amssymb}            % For additional math symbols
\usepackage{hyperref}           % For hyperlinks
\usepackage{lipsum}             % For generating dummy text

\title{Sample LaTeX Document with Title Page}
\author{First Author \\ Department of Something \\ University of Somewhere \\ \texttt{first.author@example.com}
\and Second Author \\ Department of Anything \\ University of Anywhere \\ \texttt{second.author@example.com}}
\date{\today}  % Sets date to current date

\begin{document}

% Title Page
\begin{titlepage}
    \centering
    \vspace*{2cm}  % Vertical space before title

    {\Huge Sample LaTeX Document with Title Page}\\[1.5cm]  % Title of the document

    
\end{titlepage}

\maketitle  % Generates the title

\begin{abstract}
This is the main abstract section. Here you provide a summary of the entire document. It should succinctly encapsulate the main points and findings.
\end{abstract}

\section{Introduction}
\label{sec:introduction}
This is the introduction section. In this section, you provide background information about your topic and explain the motivation behind your work. For example, you might discuss the relevance of the research question, existing literature, and any gaps that your research aims to fill.

\section{Methods}
\label{sec:methods}
In this section, you describe the methods used to conduct your research. You may include details about your experimental design, data collection procedures, and analysis methods. Use equations if necessary:

\begin{equation}
E = mc^2
\end{equation}

\section{Results}
\label{sec:results}
This section presents the results of your research. You might include tables, graphs, and other figures to illustrate your findings. 

\begin{table}[h]
\centering
\caption{Sample Table}
\begin{tabular}{|c|c|c|}
\hline
Column 1 & Column 2 & Column 3 \\
\hline
Data 1 & Data 2 & Data 3 \\
Data 4 & Data 5 & Data 6 \\
Data 7 & Data 8 & Data 9 \\
\hline
\end{tabular}
\end{table}

\section{Conclusion}
\label{sec:conclusion}
This section summarizes the main findings of your research and discusses their implications. You may also suggest directions for future research.

\begin{thebibliography}{99}
\bibitem{ref1}
Author, A. (Year). Title of the work. \textit{Journal Name}, \textit{Volume}(Issue), Page range.

\bibitem{ref2}
Author, B. (Year). Title of the work. \textit{Journal Name}, \textit{Volume}(Issue), Page range.

\end{thebibliography}

\end{document}
